\subsection{Discussion}
\label{Discussion2}
\setlength{\parindent}{20pt}
\justifying

\noindent
Experiment two provided a useful replication of experiment one, with some extra nuance present in participant responses. The addition of a morpheme spelling task prior to any training seemed to have supported morpheme learning a touch. Familiarity with the study stimuli decreased reaction times without compromising accuracy, indicating greater ease with the material. This was affected by the proportion of time participants spent speaking in a fourth language. In the absence of further information, it is unclear to the authors as to why this relationship might exist. 

The morpheme familiarity effect was also present in this version of the experiment, and this was modulated by multilingual experience. This measure increased the overall proportion of ``yes'' responses, indicating that the more experience one has with the wide range of linguistic combinations existing across languages, the more likely a new combination is judged as coherent, even if this combination does not include any previously known chunk. This was also the case with the L2 History PC component, which showed that the earlier one becomes a bilingual, the more all linguistic stimuli seem plausible. This effect was also present in participants more balanced across their first two languages. Frequent and flexible interaction of one's two main languages seem to increase their acceptance of new linguistic stimuli. Multilingual balance also played a role in participant responses: participants who scores higher on cosine similarity also were better able to differentiate the different testing conditions, and accept significantly more combinations of new and unfamiliar chunks.

Mirroring results from experiment one, experiment two also saw mean 2M responses at chance, indicating that participants on the whole were unable to pick apart congruent and incongruent 2M combinations. However, participants speaking language with multiple different types of scripts did perform above chance. Due to their experience manipulating linguistic items in different alphabets, they may have had an easier time picking up on BACS-1 stimulus characteristics. This may have more similarities to them as a foreign language script rather than something completely alien, which may have been the experience of same-script multilinguals.

This prompted further reflection on the study design. While the original stimulus set had been conceived as similar to true languages, and then converted into BACS-1 out of the necessity to erase any resemblance to any true linguistic chunk, this manipulation could be reserving more cognitive resources for the simple identification of novel letters, and take away from the processing of the morphemic chunks. Only those used to working with foreign scripts were efficient enough in single letter processing to devote enough resources to the co-occurrences within the stimuli to perform above chance. Therefore, if a version of the experiment could present the stimuli in a Latin script, this might ease all participants' cognitive load and allow for better performance. 